\chapter{マシンの夜の仕事}
% full version: chapters/20_sleep.tex

睡眠はマシンの電源を切ることではなく、モードの切り替えである。それは放送局を見るとわかる。深い睡眠では、アセチルコリン、ノルアドレナリン、セロトニンのほぼすべてが絞られる。夢を見る睡眠では、アセチルコリンが昼間のレベルに戻り、ノルアドレナリンとセロトニンはほとんど黙り、筋肉への出力はブレーキで切られる。だから夢の中で逃げている相手から、実際に走って逃げ出したりはしない。これはすべて測定されている。

\section*{いつ眠るべきか}

起きている間じゅう一日かけて充電され、夜に放電するコンデンサを想像してほしい。この電荷が「睡眠圧」である。電荷が上の閾値に達すると眠りに落ち、下の閾値まで下がると目が覚める。そして閾値そのものは、先に見た1日のリズムとともになめらかに上下する。この単純な回路はひとりでに、16時間の覚醒と8時間の睡眠に落ち着く。そして、徹夜のあとで眠る時間が2倍になるのではなく、眠りが深くなる理由も説明する。満充電のコンデンサは速く放電するのだ。「電荷」の役の候補は、脳が働くうちに蓄積する物質である。それが本当にその物質かどうかは、まだ仮説だ。

\pencilfig{120}{%
% figures/sleep_{S,H,L}.dat from models/sleep_models.py, t in hours
\begin{scope}[x=0.1cm, y=2.6cm]
\foreach \a/\b in {0/4.3,21.2/29.3,45.4/53.4,69.5/72}{\fill[pcBlue, opacity=0.1] (\a,0) rectangle (\b,1);}
\draw[pen] (0,0) -- (72,0);
\draw[penlite=pcRed, densely dashed] plot file {figures/sleep_H.dat};
\draw[penlite=pcGreen, densely dashed] plot file {figures/sleep_L.dat};
\draw[pcBlue, line width=1.0pt] plot file {figures/sleep_S.dat};
\foreach \t/\l in {0/0,24/1日,48/2日,72/3日}{\draw[pcGray, line width=0.5pt] (\t,0) -- (\t,-0.04); \node[lbl, anchor=base] at (\t,-0.16) {\l};}
\node[lbl, text=pcBlue] at (25.25,0.93) {睡眠}; \node[lbl, text=pcBlue] at (49.4,0.93) {睡眠};
\end{scope}
\node[lbl, text=pcRed, anchor=west] at (7.3,2.0) {寝つく};
\node[lbl, text=pcGreen!70!black, anchor=west] at (7.3,0.62) {目覚める};
\node[lbl, text=pcBlue, anchor=west] at (7.3,1.35) {睡眠圧};
}{睡眠圧は昼にたまり、夜に消えていく。上の閾値で眠りに落ち、下の閾値で目覚める。2つの閾値はどちらも1日の周期で揺れる。}

\section*{ゆっくりしたシーソー}

深い睡眠では、脳の領域全体が同期して、オンになったりオフになったりする。1秒に1回より遅いペースだ。これはおなじみの回路である。自分で自分を興奮させ、疲れていく群れ、つまり発振器だ。昼間はアセチルコリンがニューロンの疲労を抑えるので、同じ群れが安定して働く。夜はアセチルコリンが少なく、ネットワークは揺れ始める。

\pencilfig{20}{\pencilart{20}}{夜、ネットワークはゆっくり揺れる。ほぼすべてのニューロンが働いたかと思えば、ほぼすべてが黙る。}

\section*{早送り}

このシーソーの中で、最も興味深いことが起こる。昼間に起きたニューロンの発火の連鎖が、もう一度「再生」されるのだ。しかも何倍も速く、記憶の領域ではおよそ20倍、大脳皮質では数倍の速さで。こうした再生とゆっくりしたシーソーとの同期を人工的に強めると、記憶がよくなる。これは因果を示す実験である。何のためか。有力な仮説はこうだ。ゆっくりした長期記憶は、新しいものが以前のものを上書きしないよう、古いものと混ぜた反復から学ぶ。エンジニアはこの災難を知っている。新しい課題で訓練したニューラルネットワークは、古い例を繰り返さないと古い課題を忘れてしまう。

\section*{結合の刈り込み}

もう一つ測定されていることがある。睡眠のあと、ニューロン間の接点は平均して小さくなる。その大きさに比例して、数パーセントから数十パーセント。すべての音量を一度に少しずつ下げたようなものだ。比率は保たれ、学んだことは残り、明日の学習のための余地が空く。これが睡眠の主な役割だというのは仮説である。

\section*{夢と掃除}

夢を見る睡眠は、ベンチ試験に似ている。マシンは全力で動いているが、入力は内部の信号に置き換えられ、出力は切られている。それが何のためなのかはわからない。睡眠中に脳が老廃物を「洗い流す」という考えはまだ議論の的で、賛成する実験も反対する実験もある。

最も驚くべきこと。睡眠は局所的に起こりうる。長く起きていると、起きている人の大脳皮質の一部の領域が一瞬「眠り」、仕事の最中に黙り込むのだ。

\fullversion{20}
